\documentclass[12pt]{article}

%%%%%%%%%%%%
% packages %
%%%%%%%%%%%%
\usepackage[hmargin=0.5in,vmargin=1in]{geometry}
\usepackage{hyperref}
% \usepackage[onehalfspacing]{setspace}
\usepackage{tikz}
\usepackage{xcolor}

\definecolor{blue}{RGB}{98,119,204}
\definecolor{red}{RGB}{153,0,0}
\definecolor{brown}{RGB}{153,77,0}
\definecolor{green}{RGB}{0,153,0}
\definecolor{MyBackground}{RGB}{250,255,255}

%%%%%%%%%%%%%%%%%
% user commands %
%%%%%%%%%%%%%%%%%

%%%%%%%%%%%%
% document %
%%%%%%%%%%%%
\begin{document}

% make smaller title
\begin{center}
    {\large\bfseries Research statement, June 2026} \\
    Willard Robinson, Ph.D. student, North Carolina State Department of Economics
\end{center}

\medskip

I am an economist who studies how institutions allocate environmental and public resources, and how the frictions and information problems built into those institutions shape the outcomes they produce. I work across the methodical range of applied economics---structural welfare estimation, reduced-form causal inference, and theory---but a common theme runs throughout: when a market or mechanism is designed to deliver an efficient allocation, what happens once we take seriously the costs and behavioral responses that the idealized (``textbook'') versions assume away? My dissertation purseues this question in the context of carbon offsets and outdoor recreation, and two further projects extend it to public lands funding and to centralized assignment mechanisms.

\section{Ongoing research}

\paragraph{Job market paper: Offset markets with transaction costs}

My job market paper develops a theoretical model of offset markets that allows for non-zero transaction costs. Offsets let firms in a cap-and-trade system meet their obligations by funding abatement elsewhere, and nearly every theoretical treatment of these markets---as well as the policy solutions they imply---assumes a frictionless market. In such a setting, introducing offsets acts as a pure reduction in compliance costs. I show that this assumption holds more weight than is previously discussed.

The first result is that offset markets can generate negative welfare. Only when transaction costs and asymmetric information are combined can markets become welfare-reducing. In such cases, the optimal choice is to shut down the market.

The second result is the more surprising one, and the one I lead with: non-Pigouvian pricing is sometimes optimal. The standard prescription is to price emissions at marginal damage, but here a price set deliberately away from the Pigouvian benchmark does screening work that the ``correct'' price cannot. By exploiting the participation margin---the way fixed and proportional costs sort potential providers---a regulator can use the price itself as a screening instrument to separate additional from non-additional projects, achieving higher welfare than marginal-damage pricing delivers in a market it cannot directly monitor. The transaction cost, normally treated as pure deadweight to be minimized, becomes part of the mechanism that makes screening possible.

I lead with this paper deliberately. It is a theory contribution that the reduced-form literature cannot deliver, and it overturns a result practitioners treat as settled---that offsets weakly improve welfare and that emissions should be priced at marginal damage---with a mechanism-design argument grounded in the frictions my empirical work documents. Many of the largest offset programs in the world have high and heterogeneous participation costs and imperfectly observable counterfactuals. As a result, much of the existing literature and policy discussion has focused on improving the identification of additionality through stricter enforcement, better monitoring, and increasingly sophisticated uses of satelite imagery and machine learning. While these advances improve the regulator's ability to distinguish additional from non-additional, they do not address the underlying participation incentives which also drive non-additionality and welfare issues. In particular, costly distortions to the market mean that even perfectly measured additionality does not imply a welfare-efficient outcome. In that case, investments in better monitoring and enforcement may improve the performance of a fundamentally flawed market design rather than resolve its underlying weaknesses. By incorporating transaction costs into a model of policy-generated offset markets, I characterize when the underlying market design is sufficiently distorted that the offset program should be substantially reformed or abandoned altogether.

\paragraph{Empirical work: Transaction costs in California's offset market}

My empirical research is motivated by a pattern in the way firms act in California's cap-and-trade program. Offsets routinely trade at a substantial discount to the permits generated by the government. Yet a significant portion of firms never use them. This seemingly non-cost minimizing behavior is rationalized by the existence of transaction costs required to use offsets. Determining the success of the market rests crucially on the magnitude of these transaction costs. Using firm-level compliance data from 2013 to 2023, I try to measure these transaction costs and use these to determine whether offsets generate ``true'' economic value.

The central finding is that the binding friction is a fixed adoption cost to enerting the offset market, evidenced through the behavior of firms. California limits the total number of offsets any firm can use, and I leverage this policy peculiarity in order to determine how severe the transaction costs must be in order to rationalize the firm behavior. The intuition is this: if a firm is observed not using offsets, then its barrier to offset adoption must be at least as large as the total potential cost savings it could have earned if it used its offset limit. At the same time, a firm that is observed using offsets must have had a transaction cost no greater than the dollar cost savings it earned. I consider this leveraging of policy rules to come up with creative empirical strategies emblematic of the empirical work I intend to do in the market-based policy space.

\paragraph{Climate change-induced fish migration and angler welfare}

My third chapter turns to nonmarket valuation and the welfare consequences of climate change for outdoor recreation. As oceans warm, fish migrate toward cooler waters, so anglers in warmer regions may lose valued stocks while those further north gain them. I estimate the welfare consequences of this redistribution for shoreline anglers along the Atlantic Coast and Gulf of Mexico using a random utility travel-cost model in the tradition of recreation demand estimation. The methodological core is a chain that makes angler utility endogenous to climate: catch rates enter the site-choice and trip-frequency decisions, catch rates depend on local fish stocks, and stocks evolve under a climate-driven surplus-production model. This lets me trace an ``indirect'' welfare channel---climate acting on welfare through stocks and catch rates---that the literature has not previously quantified, and to simulate welfare gains and losses across sites under alternative warming scenarios. I have established a preliminary positive relationship between fish stocks and recreational catch rates at the state-year level; the remaining work is principally in handling the complex sampling design of the federal recreation survey microdata so that the welfare estimates can be constructed reliably at the site level.

\paragraph{Public lands funding: The Federal Land Recreation Enhancement Act}

I was invited to be a graduate fellow at the Property and Environment Research Center (PERC) in the summer of 2025 to work on this standalone project. I examine whether the 2004 Federal Land Recreation Enhancement Act, which gave national park managers greater local control over how user-fee revenue is collected and spent, crowded out central federal appropriations. To study this I digitized decades of National Park Service budget documents, assembling a panel of more than 10,000 park-year observations spanning 1986--2019, and compare fee-charging to non-fee parks before and after the reform using difference-in-differences, trend-break, and synthetic-control designs. I find little evidence of systematic crowding out, suggesting the policy added funding rather than merely redistributing it. The paper also carries a methodological lesson germane to my other work: apparent effects that are significant under park and year fixed effects disappear once park-specific time trends are included, a reminder of how much trending dependent variables can mislead. This project shares the data-construction signature of my dissertation and extends my interest in centralized-versus-local control of resource allocation to the public-lands setting.

\paragraph{Mechanism design in partially centralized matching markets}

A second standalone project, joint with Umut Dur, applies matching theory outside the environmental domain and demonstrates the breadth of my formal training. We study NC State's centralized engineering major-assignment mechanism, which closely resembles the Boston mechanism and therefore appears vulnerable to strategic misreporting. We show that under plausible assumptions about how students update their preferences over majors during a common first year---modeling preferences via a GPA-based ``match quality'' signal---the Boston mechanism and student-proposing deferred acceptance yield the same outcome, so the mechanism is in effect strategy-proof and first-choice-maximizing. Survey evidence that students overwhelmingly report truthfully is consistent with the model. The work connects to my broader interest in the design of centralized allocation mechanisms under capacity constraints and incomplete information.

\section{Agenda}

Taken together, these projects sketch a research program organized around a single tension: the gap between how institutions for allocating environmental and public resources are designed to work and how they perform once frictions and human behavior are considered more seriously. Going forward I plan to push the screening logic of the offsets model further---characterizing optimal non-Pigouvian price schedules and baseline rules under richer information structures---finish the site-level welfare estimation in the angler-welfare chapter, and extend the transaction-cost lens to other payments-for-ecosystem-services settings such as biodiversity, wetlands, and endangered-species offsets, where the same frictions arise but have been studied far less, particularly on the demand side. Methodologically, I intend to keep working at the seam between structural and reduced-form approaches, and to keep building the original datasets that can help answer these questions.


\end{document}
